\section{怎样归纳中心思想}

中心思想（在文学作品和一般记叙文中也称作主题思想）是文章的灵魂，起着统帅全文的作用。因此，在分析文章时，必须要掌握全文的中心思想。

中心思想在记叙文、议论文和说明文中表达的方法不同，因而归纳中心思想时的表达方式也不一样。

我们先谈谈怎样归纳记叙文的中心思想。

记叙文的特点，一般是通过对人物形象的具体刻画和对事件的发生、发展过程的具体描述来反映思想内容的，因而文章的中心思想往往可从两个方面来概括：一个方面是本文主要写了什么人和事；一个方面是这些人和事的描述反映了作者什么样的思想观点和感情。具体的表述方式一般是：

“本文（指明具体体裁亦可）通过对某人某事的记叙（或描写），反映了（或表现了）······的现实（或形象），歌颂了（或赞扬了）······的品质（或精神），暴露了（或揭示了）······的罪行（或问题、实质），表达了（或流露了）作者······的思想感情。”

其中第一句话是必不可少的，因为它是作者表达思想感情的基础和前提；以下的几句话则应根据具体文章的思想内容来决定取舍，不一定每篇记叙文的中心思想都要写这么多项目。在语言表达上，第一句话要尽量概括、简练，只要把记叙的人和事点到即可。问题是下面的话该怎样归纳出来呢？这可从以下几方面入手。

1. 分析人物和事件所反映的思想意义。我们只要弄清作者所写的人物是什么样的人物，作者对他的态度怎样，所写的事件说明什么问题，宣传什么道理，就能从中概括出中心思想。例如《守财奴》（部编语文高中第三册），作者所写的葛朗台是一个为夺取别人财产而施展抢和骗的手段，达到爱财如命、贪婪成癖地步的资产阶级暴发户、投机商，通过他妄图夺取金梳妆匣、剥夺女儿财产继承权和临终抓取镀金十字架等一系列事件的描写，深刻地揭露了资本主义社会人与人之间赤裸裸的金钱关系。我们只要分析了文章中人物和事件所反映的思想，并将它们综合起来，就可以归纳出全文的中心思想了。再如鲁迅小说《祝福》（部编语文高中第二册），写了旧中国农村里一个饱受压迫和侮辱的劳动妇女祥林嫂终于在封建礼教、鬼神迷信和宗法制度的摧残下悲惨死去的故事。作者通过对以鲁四老爷为代表的杀死千百万祥林嫂的社会环境的揭露，深刻显示了那个“病态社会”中压迫和摧残妇女的“畸形道德”的虚伪性和残酷性，指明了统治着那个社会的封建礼教和迷信思想的吃人本质，启示我们如果劳动人民不摆脱封建礼教与迷信思想的束缚，就不能避免悲剧的命运。经过这样分析之后，这篇小说的中心思想也就非常清楚了。

2. 根据段意归纳中心思想。这是一种由分散到集中、由部分到全体的归纳方法，我们只要把全篇各个意义段落集中起来，逐层统一，就能提炼出全文的中心思想。例如莫泊桑小说《项链》（部编语文高中第四册），全文共五个部分，分别写了：小资产阶级出身的路瓦裁夫人羡慕资产阶级的奢华生活，借项链的虚荣；丢项链的惊慌；赔项链的劳
苦；发现是假项链后的懊恼。我们可以在弄清小说基本情节和作者所持态度的基础上，把全文的中心思想归纳出来：小说写了路瓦栽夫人在一次晚会中丢失了一挂借来的项链，逼得她用十年时间历尽艰辛攒钱赔偿的故事，反映了一个爱虚荣的小资产阶级妇女的不幸遭遇，嘲讽了小资产阶级的虚荣心理，揭露了资产阶级思想对人的毒害。再如鲁迅小说《药》（部编语文高中第一册），全文共四个部分，分别写了：贫民华老栓为替儿子治病而去刑场买“药”——蘸着烈士夏瑜鲜血的人血馒头；华老栓夫妇照管儿子华小栓稀里胡涂地吃“药”；借满脸横肉的康大叔谈“药”，写出为推翻封建制度而英勇献身的革命者夏瑜不被群众理解的悲剧；华大妈和夏四奶奶同为儿子上坟而相遇，却全不了解夏瑜就义的意义。我们在弄清小说的基本情节以及作者给这些情节所赋予的内在含义和悲剧色彩之后，可以把全文的中心思想归纳出来：作品通过革命者夏瑜的鲜血竟成了贫民华老栓夫妇为儿子治病的“药”这一悲剧事件，不仅揭露了封建统治阶级镇压革命和愚弄人民的反动本质，更批判了资产阶级旧民主主义革命严重脱离群众的错误，启示人们只有发动群众，革命才能胜利，华、夏才能得救。

3.注意文章的题目、开头和结尾。有的文章题目本身就揭示了中心思想或者暗示了中心思想。例如初中语文课本中的《生的伟大，死的光荣》、《大渡河畔英雄多》、《接过前辈的扁担》、《鞠躬尽瘁》等题目，就是这样。即使有些文章的题目本身并不能揭示中心思想，它们也从不同方面提示了文章中心之所在。例如高中语文课本中的《为了六十一个阶级弟兄》、《樱花赞》、《爱国学者顾炎武》、《装在套子里的人》等题目，就是这样。当然，更多的文章题目只是体现出文章中所写的人物、主要事件、描写对象等，例如高中语文课本中的《范爱农》、《包身工》、《“老虎团”的结局》、《山地回忆》、《火刑》、《秋色赋》等篇，就是这样，我们可以围绕题目，发掘作者的爱憎感情和思想倾向，自能给我们归纳中心思想指出路子来。

另外，有的文章在开头或结尾直接阐明了中心思想。如魏巍《谁是最可爱的人》的首段和尾段，碧野的《天山景物记》的首段和尾段，唐弢《琐忆》的开头，吴晗《爱国学者顾炎武》的开头，夏衍《包身工》的结尾，鲁迅《一件小事》的结尾，等等，都明确地写出了文章的主旨。

4.注意文章的议论和抒情成分。记叙文中的议论和抒情成分，往往直接表达了作者对所记叙的人和事的看法，同时揭示出事物的本质，表达了中心思想。如杨朔的《海市》、郑文光的《火刑》、茅盾的《风景谈》、方纪的《挥手之间》，以及《谁是最可爱的人》、《天山景物记》、《一件小事》、《荔枝蜜》等文中的议论和抒情成分就包含了全篇的中心思想。

我们只要综合运用这些有效手段，善于抓住文章思想内容的主要方面，并能够简要地表达出来，就行了。

下面我们谈谈怎样归纳议论文的中心思想。

议论文的特点是直接说理，作者通过对某一问题的提出与分析来阐明他所持的见解和主张，作者的见解、主张也就是文章的中心论点。因此，议论文的中心思想要在把握文章的中心论点和论证过程的基础上加以概括，因为中心思想的核心是中心论点，还要概括文章的主要内容以及作者提出的希望、要求、号召等内容。议论文中心思想具体表述的方式，一般是：

“本文分析了（或总结了）······的问题（或经验），阐明了······的道理，肯定了······的看法，谴责了（或驳斥了、揭露了）……的罪行（或谬论、嘴脸），指出了（或指明了）······的道路（或方向），号召（或鼓舞）······”

当然，归纳一篇议论文的中心思想，不必把上述的项目一一列举，而要根据具体内容进行取舍。例如《改造我们的学习》（高中第三册语文）的中心论点就是开头的一句话：“我主张将我们全党的学习方法和学习制度改造一下。”然后分四个部分来写。前三部分反复论述为什么要改造我们的学习，第四部分说明怎样改造我们的学习。我们可以根据全文论述的问题，归纳其中心思想如下：“本文总结了在学风问题上的历史经验和教训，深刻批判了理论脱离实际的反马克思主义的学风，指出了马列主义的普遍真理和中国革命的具体实践相结合是我党唯一正确的学风，阐明了必须改造全党的学习方法和学习制度的道理，并为如何改造我们的学习指明了方向。”

在归纳中心思想时，应注意如下几点：

\textbf{1. 要找准文章的中心论点。}

不少议论文的中心论点就表现在文章的题目上，例如《反对自由主义》、《实践是检验真理的标准》、《作一个好的党员，建设一个好的党》、《学习白求恩》，等等。有的在文章开头直截了当地提出了中心论点，如《改造我们的学习》、《放下包袱、开动机器》（初中第六册语文）、《谈骨气》（初中第一册语文）等。有的在文章结尾作结论时，才明确中心论点，如《民族的科学的大众的文化》（高中第二册语文）、《拿来主义》（高中第二册语文）、《什么是知识》（初中第五册语文）等。

\textbf{2. 根据段意归纳中心思想。}

这和记叙文一样，在划分段落、写出段意的基础上，我们把各个意义段的段意进行概括，逐层统一，最后确定中心思想。我们刚才对《改造我们的学习》的中心思想的归纳就是这样做的。

再如鲁迅的著名杂文《拿来主义》（高中第二册语文）。第一意义段（1—5自然段）先论述“闭关主义”和“送去主义”的错误，对比鲜明地提出了“拿来主义”；第二意义段（6—9自然段）分析“拿来”与“送来”的本质区别，运用比喻论证说明应该“要运用脑髓，放出眼光，自己来拿”；第三意义段（第10自然段）作出结论，指明正确对待文化遗产的原则。我们据此可以归纳出本文的中心思想：本文批判了国民党反动派的卖国主义和各种错误倾向，深刻而生动地阐明了马克思主义关于批判继承文化遗产的基本原则。

\textbf{3. 从整个论述过程入手，进行分析和归纳。}

有的议论文，表面找不到点明中心论点的语句，论点在论证中体现出来，要我们自己去领悟概括。我们可以分析全文的论述过程，弄清文章的逻辑结构和内在联系，看看文章究竟是提出了什么问题，又是怎样分析和解决的，然后把它们综合起来，也就明确了全文的中心思想。在这类文章中，总有概括性的语句，我们可以把它作为归纳中心思想的一种依据。例如鲁迅的著名散文《论雷峰塔的倒掉》（初中第四册语文），全文没有点明中心论点的句子，但是论述的过程是很清楚的：先从“听说，杭州西湖上的雷峰塔倒掉了”谈起，然后回忆幼时所听的关于白蛇与许仙结婚、终于被法海和尚镇在雷峰塔下的故事，“那时我唯一的希望，就在这雷峰塔倒掉”，接着说到“现在，他居然倒掉了，则普天之下的人民，其欣喜为何如？”并引事实作证，结尾写雷峰塔终于倒掉了，而法海却变成螳螂和尚，永世躲在蟹壳里，“活该”！从上述的论证和概括性的句子来看，本文的中心思想就是：深刻地阐明了“一切压在人民头上的‘雷峰塔’终将倒掉的道理，强烈地表达了对以白蛇为代表的追求幸福生活的人民的同情，对以法海为代表的压迫者的憎恨的思想感情。

最后我们来谈谈怎样归纳说明文的中心思想。

说明文的基本要求是说明事物的特征和本质，因此，用有关的材料写成文章，说明事物的特征和本质，使人增长知识，受到教育，也就表达了文章的中心思想。归纳说明文的中心思想的具体表达方式，一般是：

“本文说明了······的特征（形状、构造、特性、功能、用途、过程等），指出了了······的本质（历史意义、认识价值等）。”

例如《中国石拱桥》一文，先说明中国石拱桥的历史特别悠久，然后以赵州桥和芦沟桥为例，说明我国石拱桥在设计和施工上的独特创造与高度的技术水平以及获得成就的原因，最后说明我国解放后石拱桥建筑技术的新发展和新成就。我们对按照各段说明的内容加以归纳，说明文的中心思想就显示出来了。《中国石拱桥》中心思想可以概括如下：

“本文说明了我国石拱桥的悠久历史和建筑上的杰出成就，指出了在优越的社会主义制度下我国桥梁事业飞跃发展的现实和光明前景。”

“说明了了······”一句概括了事物的特征，“指出了了······”一句揭示了事物的本质。一般说明文的中心思想都可按照这种方法加以概述。
\newpage